\section{Post-Newtonian and solar-system constraints (Einstein gravity with an information scalar)}
\label{app:ppn_solar_system}

This appendix positions the CAP-II source sector relative to standard scalar--tensor frameworks and records the minimal post-Newtonian (PPN) implications needed for solar-system tests.
The goal is not to re-derive the full PPN formalism, but to make explicit what is fixed by the Einsteinian geometric skeleton and what must be constrained by the information-sector constitutive inputs.

\subsection{Relation to scalar--tensor theories}
\label{app:ppn_relation_scalar_tensor}

The minimal CAP-II closure uses the action \eqref{eq:omega_action_cap2}, i.e.\ Einstein--Hilbert gravity with an additional scalar field $\chi$ and covariant matter $\mathcal{L}_m(g,\psi_m)$.
In particular, in this manuscript:
\begin{itemize}
  \item the gravitational coupling $G$ is constant in the action, and
  \item matter couples minimally to the metric (no direct $\chi$--matter coupling is assumed in $\mathcal{L}_m$).
\end{itemize}
This differs from Brans--Dicke and more general scalar--tensor theories, where the scalar typically modifies the effective gravitational coupling (Jordan-frame $F(\varphi)R$ terms) and/or couples nontrivially to matter in the Einstein frame; see, e.g., \cite{BransDicke1961,Will1993,Will2014}.
\par
\noindent Consequently, CAP-II in its minimal form is a metric theory with Einstein field equations and an additional stress tensor contribution.
Deviations from Schwarzschild in exterior regions arise not from a modified left-hand side, but from nontrivial information-sector stress-energy $T^{\mathrm{info}}_{\mu\nu}$ (and from any additional matter sector used in a given realization).

\subsection{PPN parameters in the minimal closure}
\label{app:ppn_parameters}

In standard PPN gauge (isotropic spatial coordinates), the weak-field metric of a static source is written schematically as
\begin{align}
  g_{00} &= -1 + 2U - 2\beta U^2 + O(U^3),\\
  g_{ij} &= \left(1 + 2\gamma U + O(U^2)\right)\delta_{ij},
\end{align}
where $U=GM/r$ is the Newtonian potential and $(\gamma,\beta)$ are theory parameters; see, e.g., \cite{Will1993,Will2014}.
For Einstein gravity with constant $G$ and minimally coupled matter, one has
\begin{equation}
  \gamma = 1,
  \qquad
  \beta = 1,
\end{equation}
independent of the detailed constitution of the source.
In CAP-II this statement applies to the \emph{geometric skeleton}: the post-Newtonian coefficients are those of GR in the regime where the exterior metric is well-approximated by a Schwarzschild-like solution.
\par
\noindent Empirically, solar-system tracking constrains $\gamma$ at the level $|\gamma-1|\ll 1$; see \cite{Will2014} for current bounds and a consolidated discussion.

\paragraph{A check in a non-vacuum exterior (massless scalar/JNW family).}
To make the above concrete in a setting where the exterior is not vacuum, consider Einstein gravity coupled to a free massless scalar (the $V\equiv 0$ limit of \eqref{eq:omega_action_cap2} with $\lambda_F>0$).
The generic static spherically symmetric solution is the Fisher/JNW family \cite{Fisher1948,JanisNewmanWinicour1968}, which can be written in curvature coordinates as
\begin{equation}
  \dd s^2
  = -\left(1-\frac{b}{r}\right)^{\nu}\dd t^2
  + \left(1-\frac{b}{r}\right)^{-\nu}\dd r^2
  + \left(1-\frac{b}{r}\right)^{1-\nu} r^2\dd\Omega^2,
  \qquad 0<\nu\le 1,
  \label{eq:jnw_metric}
\end{equation}
with $\nu=1$ recovering Schwarzschild.
Transforming to an isotropic radius $\rho$ and expanding at large $\rho$ yields
\begin{align}
  g_{00} &= -1 + \frac{\nu b}{\rho} - \frac{\nu^2 b^2}{2\rho^2} + O(\rho^{-3}),\\
  g_{ij} &= \left(1 + \frac{\nu b}{\rho} + O(\rho^{-2})\right)\delta_{ij},
\end{align}
so identifying $GM=\nu b/2$ one reads off $\gamma=\beta=1$ at post-Newtonian order.
Thus, even when the exterior contains scalar stress-energy, the leading PPN coefficients remain Einsteinian; constraints arise from higher-order terms and from departures from a Schwarzschild benchmark in a chosen exterior window.

\subsection{Standard solar-system observables (leading order)}
\label{app:ppn_observables}

At leading post-Newtonian order, the classic tests depend only on $\gamma$ and $\beta$.
In particular, for light deflection at impact parameter $b$ and Shapiro time delay for a radar signal with endpoints at radii $r_1,r_2$, one has (restoring $c$ for clarity)
\begin{align}
  \delta\theta_{\mathrm{light}}
  &= \frac{1+\gamma}{2}\,\frac{4GM}{b c^2}
  = \frac{4GM}{b c^2}
  \qquad (\gamma=1),\\
  \Delta t_{\mathrm{Shapiro}}
  &\approx \frac{1+\gamma}{2}\,\frac{4GM}{c^3}\log\!\left(\frac{4r_1 r_2}{b^2}\right)
  = \frac{2GM}{c^3}\log\!\left(\frac{4r_1 r_2}{b^2}\right)
  \qquad (\gamma=1),
\end{align}
see \cite{Shapiro1964,Will1993,Will2014} for conventions and refinements.
\par
\noindent In CAP-II's operational dictionary (Section~\ref{subsec:lapse_dictionary}), the weak-field potential is $\Phi=-\log\lapse=\log(\kappa/\kappa_0)$ and $g_{00}\approx-\lapse^2$ in a shiftless gauge.
Therefore, once a dataset provides $\kappa(x)$ (by compilation logs or by the WS interface) and a benchmark exterior interpretation identifies $GM$ in an exterior window (Section~\ref{subsec:schwarzschild_fit}), the standard leading-order solar-system observables follow in that same window.

\subsection{Constraints on the information scalar: massive vs.\ light regimes}
\label{app:ppn_constraints_on_scalar}

The parameter dependence specific to CAP-II enters through the source sector $T^{\mathrm{info}}_{\mu\nu}$ and the $\kappa\mapsto\chi$ constitutive map.
Solar-system consistency is therefore most naturally phrased as a constraint that the information sector does not generate observable departures from the benchmark exterior templates in the relevant window.

\paragraph{Massive (pinned) regime.}
Assume $V$ has a local minimum at the background $\chi=\chi_0$ with $V'(\chi_0)=0$.
Linearizing \eqref{eq:chi_eom_cap2} gives
\begin{equation}
  \square\,\delta\chi - m_\chi^2\,\delta\chi = 0,
  \qquad
  m_\chi^2 := \frac{V''(\chi_0)}{2\lambda_F},
\end{equation}
so $\delta\chi$ is Yukawa-suppressed beyond the Compton length $\ell_\chi:=m_\chi^{-1}$.
In this regime, an exterior window with length scales $r\gg \ell_\chi$ is effectively Schwarzschild-like up to exponentially small corrections from the scalar gradients, and the benchmark fits of Section~\ref{subsec:schwarzschild_fit} are self-consistent.
\par
\noindent A conservative solar-system consistency condition is that the scalar be short-ranged compared to the smallest impact parameters $b$ used in classic time-delay/deflection tests, i.e.
\begin{equation}
  m_\chi\,b \gg 1
  \qquad\text{(or equivalently } \ell_\chi \ll b\text{)}.
\end{equation}
For example, taking $b$ of order a solar radius provides a concrete benchmark scale.
In terms of the $\rho$-expansion family \eqref{eq:V_rho_expansion} with $V(\chi^2)=\widehat V(\rho)$ and $\rho=\chi^2$, one has near the minimum $\rho_0=\chi_0^2$ that $V''(\chi_0)=4\chi_0^2\,\widehat V''(\rho_0)$, so
\begin{equation}
  m_\chi^2 = \frac{2\chi_0^2}{\lambda_F}\,\widehat V''(\rho_0).
\end{equation}
Thus a lower bound on $\widehat V''(\rho_0)$ (given $\lambda_F$ and $\chi_0$) is a directly interpretable short-range condition.
\par
\noindent In the explicit quadratic family \eqref{eq:V_rho_expansion}, $\widehat V''(\rho_0)=m_\rho^2$, hence
\begin{equation}
  m_\chi^2 = \frac{2\rho_0}{\lambda_F}\,m_\rho^2,
  \qquad (\rho_0=\chi_0^2).
\end{equation}
\par
\noindent The exponent $p$ enters these constraints through the constitutive identification $\chi=\sqrt{\rho_0}(\kappa/\kappa_0)^{p/2}$: in the weak-field regime $u=\log(\kappa/\kappa_0)$ is small and $\delta\chi/\chi_0 \approx (p/2)\,u$.
Thus, for a fixed operational lapse profile $u(x)$, smaller $p$ suppresses the amplitude of scalar-sector excursions and therefore the size of $T^{\mathrm{info}}_{\mu\nu}$ corrections in an exterior window.

\paragraph{Light (nearly massless) regime and JNW deviations.}
If $m_\chi$ is very small on solar-system scales, static spherically symmetric exteriors generically carry scalar hair and are described by Fisher/JNW-type families (or their potential-deformed analogues).
Although $\gamma=\beta=1$ at leading PPN order (Section~\ref{app:ppn_parameters}), the exact exterior metric can differ from Schwarzschild at higher orders and in the near field.
In CAP-II this regime is treated as an \emph{observable deviation channel}: departures from the Schwarzschild benchmark constrain the potential family $V(\chi^2)$ and the effective coupling scale set by $(p,\rho_0,\lambda_F)$ through the $\kappa\mapsto\chi$ identification.

\subsection{Fifth-force interpretation and screening mechanism}
\label{app:fifth_force_screening}

In the minimal closure \eqref{eq:omega_action_cap2} with minimally coupled matter, test bodies follow metric geodesics, so there is no additional \emph{direct} fifth force from a matter--scalar coupling.
Any observable deviation from GR arises indirectly through the scalar contribution to the metric sourced by $T^{\mathrm{info}}_{\mu\nu}$.
\par
\noindent Accordingly, the operative ``screening'' mechanism in this setup is the massive/pinned regime of Section~\ref{app:ppn_constraints_on_scalar}: if $m_\chi$ is large compared to the inverse length scales probed by an experiment, scalar gradients and their stress-energy are exponentially suppressed outside the source region, making exterior observables indistinguishable from GR within stated precision.
Chameleon-like screening requires an environment-dependent effective potential (typically via explicit matter coupling), which is not assumed in the minimal CAP-II closure and would constitute an additional modeling layer.

\subsection{Binary-pulsar and radiative constraints (qualitative placement)}
\label{app:binary_pulsar_constraints}

Binary-pulsar tests constrain departures from GR through strong-field dynamics and radiative channels.
In scalar--tensor theories with matter coupling, light scalars typically generate dipole radiation and are therefore tightly constrained; see \cite{Will2014}.
CAP-II in its minimal form does not introduce a direct matter--scalar coupling, so the leading dipole-radiation mechanism of Jordan-frame scalar--tensor theories is not automatically present.
Operationally, the relevant CAP-II constraint is again that the information sector remain short-ranged or sufficiently weak in the exterior/radiative zone so that waveform phase evolution and timing observables are consistent with the GR templates used in the analysis.


